\documentclass[11pt,letterpaper]{article}
\usepackage[margin=0.91in,top=0.82in,bottom=0.85in,headheight=15pt,headsep=12pt]{geometry}
\usepackage{amsmath,amssymb,amsthm,booktabs,array,xcolor,fancyhdr}
\usepackage[colorlinks=true,linkcolor=ink,citecolor=ink,urlcolor=teal]{hyperref}
\definecolor{ink}{RGB}{21,54,76}
\definecolor{teal}{RGB}{0,101,102}
\definecolor{muted}{RGB}{75,87,94}
\pagestyle{fancy}\fancyhf{}
\lhead{\small\sffamily\color{muted} Li coefficients and zero separation}
\rhead{\small\sffamily\color{muted} Research note}
\rfoot{\small\sffamily\thepage}
\renewcommand{\headrulewidth}{0.3pt}
\linespread{1.06}\setlength{\parindent}{0pt}\setlength{\parskip}{5pt plus 1pt minus 1pt}
\setlength{\abovedisplayskip}{8pt plus 2pt minus 2pt}
\setlength{\belowdisplayskip}{8pt plus 2pt minus 2pt}
\newcommand{\HH}[1]{\section*{\color{ink}#1}}
\newcommand{\subt}[1]{\par\smallskip{\sffamily\bfseries\color{teal}#1}\par\nobreak}
\newcommand{\wt}{w_a}
\newcommand{\Gal}{\mathcal G}
\DeclareMathOperator{\Rea}{Re}
\newcommand{\Z}{\mathbb Z}
\newcommand{\R}{\mathbb R}
\newcommand{\C}{\mathbb C}
\newcommand{\op}[1]{\operatorname{#1}}
\newcommand{\dd}{\mathrm d}
\newcommand{\TT}{\mathbb T}
\newcommand{\code}[1]{\texttt{#1}}
\hypersetup{pdftitle={Detecting prescribed gaps between Dirichlet L-function zeros from finite Li data},pdfauthor={OpenAI ChatGPT}}
\begin{document}
\thispagestyle{plain}
{\sffamily\small\color{teal}\bfseries ANALYTIC NUMBER THEORY}\hfill
{\small\sffamily 26 November 2026}
\vspace{7pt}

{\LARGE\bfseries\color{ink} Detecting prescribed gaps between Dirichlet\par
\vspace{1pt} \(L\)-function zeros from finite Li data\par}
\vspace{5pt}
Let $\beta(s)=L(s,\chi_{-4})$ be the Dirichlet beta function. If $\gamma$ runs over the
ordinates of the nontrivial zeros of $\zeta$ and $\eta$ over those of $\beta$, we certify
\[
  \boxed{\quad 0.361 <
  \min_{|\gamma|,|\eta|\le 22}|\gamma-\eta| < 0.429 .\quad}
\]
For the lower bound, the only function data are finite Taylor expansions of
the two completed functions at the same real point. For the upper bound,
four rigorous sign evaluations exhibit a pair. The general theorem explains
when a finite calculation of the first kind must succeed.

\subt{The finite test}
Fix primitive Dirichlet $L$-functions $A,B$, including $\zeta$, and real centers
$c_A,c_B$. Fix a window radius $T\ge0$ and a scale $a>0$. The question is whether
there are zeros with ordinates $\gamma,\eta$ such that
\[
 \gamma-c_A=\eta-c_B,\qquad |\gamma-c_A|\le T .
\]
For integers $N\ge2$ and $r\ge1$, we construct a real number $F_{r,N}$
from the first $rN$ shifted Li coefficients of each completion. Put
$W=4a^2/(a^2+T^2)$. We prove the equivalence
\[
 \text{the window contains no such pair}
 \quad\Longleftrightarrow\quad
 F_{r,N}<W^{2r}\ \text{for some }r,N .
\]
In a window free of matching pairs, every schedule with $r,N\to\infty$
eventually passes. With $A=B$ and equal centers, a diagonal subtraction gives
the corresponding test for multiple zeros. A quantitative version bounds the
distance from the prescribed gap even when the ordinates differ.

\subt{What comes from the two sources}
The readable Lean edition gives nonvanishing to the right of $1/2$ and
critical-strip zero location for Dirichlet characters, including imprimitive
ones \cite{lean}. Its reading guide traces the proof through short integral forms,
spanning, and a lower bound on integral forms. Their comparison bounds the same index-corrected covolume from both sides.
The main term cancels, bounding a prime remainder smoothed with $1/n-1/x$.
The one-sided zero criterion then gives the zero-free half-plane.
We use the resulting real ordinates. The requested article's discussion of
Li's positivity criterion suggests encoding these ordinates by a sequence
\cite{wiki,bl}. The additional operation here compares the sequences of
\emph{two} functions. Lemma~1 identifies the particular finite differences
needed for this comparison. In the pinned source,
\code{LFunction\_ne\_zero\_of\_half\_lt\_re} excludes the principal pole
and \code{LFunction\_eq\_zero\_re} locates the zeros in the open critical strip.

\newpage
\HH{1. From a Taylor expansion to a circle measure}
Write $H$ for the entire completion of a primitive Dirichlet $L$-function,
with the polar factors removed in conductor one. For zeta we use\[\xi(s)=\tfrac12 s(s-1)\pi^{-s/2}\Gamma(s/2)\zeta(s).\]
These functions have order one. Their completed zeros are the nontrivial
zeros, counted with multiplicity. Their functional equation has the form
$H(s)=\epsilon\overline{H(1-\bar s)}$ with $|\epsilon|=1$ \cite{lagarias}.
The critical-line input from \cite{lean} writes each zero as
$\rho=1/2+i\gamma$, $\gamma\in\R$. For a center $c$ define
\[
 s_\pm=\tfrac12+ic\pm a,\quad d=2a,\quad S(u)=s_++\frac{du}{1-u},
 \qquad z_\gamma=\frac{\rho-s_-}{\rho-s_+},\quad
 w_\gamma=|1-z_\gamma|^2=\frac{d^2}{a^2+(\gamma-c)^2}.
\]
We have $|z_\gamma|=1$, $0<w_\gamma\le4$ and $\sum_\gamma w_\gamma<\infty$.
The last fact follows from the order-one zero count, or the convergence
of the genus-one product. In all sums a separate label represents each
occurrence of a zero. Write $[u^n]$ for coefficient extraction. Set
\[
 \ell_n=n[u^n]\log H(S(u))\ (n>0),\qquad
 \ell_0=0,\qquad \ell_{-n}=\overline{\ell_n}.
\]
These coefficients are independent of the local logarithm branch. At
$a=1/2,c=0$ this is the usual right-center Li expansion. Real shifts
of the evaluation point have a precedent in Sekatskii \cite{sek}.
For nonreal characters or nonzero $c$, the signed indices matter.

\textbf{Lemma 1.} Let $(Db)_n=2b_n-b_{n-1}-b_{n+1}$. For all $r\ge1$, $n\in\Z$,
\begin{equation}
 b_{r,n}:=-(D^r\ell)_n=\sum_\gamma w_\gamma^r z_\gamma^n .
 \label{eq:moment}
\end{equation}
\emph{Proof.} Write $G=H'/H$. The genus-one logarithmic derivative
difference cancels the exponential factor in the Hadamard product, giving
\[
 d\{G(S(u))-G(s_-)\}
 =d\sum_\rho\left(\frac1{S(u)-\rho}-\frac1{s_--\rho}\right)
 =\sum_\gamma\frac{w_\gamma}{1-z_\gamma u},\qquad |u|<1 .
\]
Indeed $S(u)-s_-=d/(1-u)$ and
$w_\gamma=-d^2/[(\rho-s_-)(\rho-s_+)]$. Both sums converge locally uniformly
by the summability above. Differentiating the defining Taylor series gives
$\sum_{n\ge1}\ell_n u^n=duG(S(u))/(1-u)^2$.
For $n\ge1$, the coefficient of $u^n$ in $dG(S(u))$ is therefore
$\ell_{n+1}-2\ell_n+\ell_{n-1}$. At zero the left side of the difference
identity equals $2\Rea\ell_1$, since $G(s_-)=-\overline{G(s_+)}$.
Conjugation handles negative indices. Finally $D$ multiplies the mode $z^n$
by $2-z-z^{-1}=|1-z|^2$. This proves \eqref{eq:moment}. \hfill$\square$

\newpage
\HH{2. A finite expression for matched ordinates}
On the unit circle put
\[
 K_N(v)=\left|\frac1N\sum_{j=0}^{N-1}v^j\right|^2,\qquad
 K_N(v)^r=\sum_{|k|\le r(N-1)}h_kv^k .
\]
The rational coefficients $h_k$ are symmetric and nonnegative because they come
from the $r$-fold convolution of $(N-|k|)/N^2$ for $|k|<N$.
Also $K_N(1)=1$, $0\le K_N\le1$, and $K_N(v)\to0$ for $v\ne1$.
For $A,B$ use their respective centers in Lemma~1, with the same scale $a$.
Define the mixed test by
\begin{equation}
 F_{r,N}=\sum_{|k|\le r(N-1)}h_k b^A_{r,k}\overline{b^B_{r,k}} .
 \label{eq:F}
\end{equation}
For the multiplicity test take $A=B$, $c_A=c_B$, and subtract $b^A_{2r,0}$
from the right side. A prescribed nonzero shift within a single function uses
the mixed test with two centers. Both versions use only indices $|n|\le rN$,
as $2r\le rN$. The formula is quadratic in the finite lists, with a linear
correction in the multiplicity test.

\textbf{Lemma 2.} Write $i,j$ for zero labels of the two functions, $z_i,v_j$
for their circle points, $w_i,w'_j$ for their weights, and put
\[
 p_{ij}=w_iw'_j,\qquad q_{ij,N}=p_{ij}K_N(z_i\bar v_j).
\]
Use all ordered pairs for the mixed test. In the multiplicity test omit
exactly the pairs of a label with itself. With these conventions,
\begin{equation}
                  F_{r,N}=\sum_{i,j}q_{ij,N}^{\,r}\ge0.
 \label{eq:pairs}
\end{equation}
\emph{Proof.} Expand the polynomial in \eqref{eq:F} and apply \eqref{eq:moment}.
All sums are absolutely convergent since $0\le q_{ij,N}\le p_{ij}\le16$ and
$C:=\sum_{i,j}p_{ij}\le(\sum_iw_i)(\sum_jw'_j)<\infty$.
In the multiplicity test each omitted diagonal contributed $w_i^{2r}$.
Two distinct labels of the same zero survive the subtraction. \hfill$\square$

\subt{What a fixed angular resolution preserves}
Put $x_i=\gamma_i-c_A$, $y_j=\eta_j-c_B$. The map
$x\mapsto (ix+a)/(ix-a)$ is injective. Thus $z_i=v_j$ exactly when $x_i=y_j$.
For every $t>0$ there are only finitely many pairs with $p_{ij}\ge t$,
since both weights are at least $t/4$. Among these, call a pair a
\emph{match} when $x_i=y_j$. As $N$ increases, all unequal pairs in this
finite set have $q_{ij,N}\to0$, whereas matched pairs have $q_{ij,N}=p_{ij}$
for every $N$. Pairs outside the set already lie below $t$. Consequently,
for all sufficiently large $N$, the full multiset of $q_{ij,N}\ge t$
equals the multiset of matched products $p_{ij}\ge t$, with multiplicities.

\newpage
\HH{3. Completeness and the first match}
\textbf{Theorem 3.} In the mixed test, there is no match with $|x_i|\le T$
if and only if $F_{r,N}<W^{2r}$ for some $r\ge1,N\ge2$, with
$W=d^2/(a^2+T^2)$. In the multiplicity test the same equivalence holds
with the threshold $2W^{2r}$. On a clear window the relevant strict
inequality eventually holds along every joint sequence $r,N\to\infty$.

\emph{Proof.} A match in the window has $q_{ij,N}=w_i^2\ge W^2$.
A repeated zero of multiplicity $m$ supplies $m(m-1)\ge2$ ordered pairs.
This proves the necessity of a clear window for a strict certificate.

Conversely, let $R$ be the supremum of matched products, with $R=0$ for an
empty set. Every positive supremum is attained, by finiteness above each
positive threshold. If the window is clear, $R<W^2$, including when a match
could have fallen exactly at the boundary. Choose $R<t<W^2$. The
finite-set argument on the preceding page gives $q_{ij,N}<t$ for all pairs
and all sufficiently large $N$. Uniformly in these $N$,
\[
 0\le \frac{F_{r,N}}{W^{2r}}
 \le \frac{C}{W^2}\left(\frac t{W^2}\right)^{r-1}\longrightarrow0 .
\]
This proves sufficiency and the assertion about joint sequences. \hfill$\square$

\textbf{Corollary 4.} Along every such joint sequence,
$F_{r,N}^{1/r}\to R$. If $R>0$, let $M$ be the number of admissible ordered
pairs realizing the largest matched product. Then there exist $b<1$ and
$N_0$ for which
\[
             0\le F_{r,N}/R^r-M\le (C/R)b^{r-1}\qquad(N\ge N_0).
\]
In particular the least absolute centered ordinate at a match equals
$\sqrt{d^2/\sqrt R-a^2}$, interpreted as $+\infty$ for $R=0$.

\emph{Proof.} When $R=0$, take any $t>0$ in the preceding estimate
$F_{r,N}\le Ct^{r-1}$. When $R>0$, remove all maximizing matched pairs.
The remaining matched supremum is strictly smaller than $R$ by finiteness
above $R/2$. Pick $b\in(0,1)$ strictly between its ratio to $R$ and $1$.
The stabilization result with threshold $bR$ gives the bound on the remainder.
The formula for the ordinate inverts $R=\wt(x)^2$, where
$\wt(x)=d^2/(a^2+x^2)$. \hfill$\square$

There is a monotone choice of approximations. The factorization
$K_{mN}(v)=K_N(v)K_m(v^N)$ implies $F_{r,mN}\le F_{r,N}$.
For fixed $N$, $F_{r,N}^{1/r}$ is the $\ell^r$ norm of a sequence over
individual pairs, so it decreases with $r$. Taking $N=2^j,r=j$ therefore
approaches $R$ from above. The theorem supplies no uniform bound on how
large $j$ must be for an arbitrary pair of functions or window.

\newpage
\HH{4. Turning a strict test into a gap bound}
Let $x_i,y_j\in[-T,T]$ denote the centered ordinates of an admissible pair.
Suppose a known symmetry gives at least $k$ pair contributions of equal value
in \eqref{eq:pairs}, including the chosen pair. One can always use $k=1$.
In the multiplicity test, reversal gives $k=2$. In our mixed example,
both functions are real self-dual and centered at zero. Simultaneous
negation gives $k=2$ since $\zeta(1/2)\ne0$, as follows from positivity of the alternating eta
series at $1/2$.

\textbf{Theorem 5.} Write $P_0=W^2$ and $Q=(F_{r,N}/k)^{1/r}$.
If $Q<P_0$, then every such pair satisfies
\begin{equation}
 |x_i-y_j|\ \ge\ d\sqrt{\frac{12}{N^2-1}
   \min\!\left\{\frac{P_0-Q}{P_0^2},\frac{16-Q}{256}\right\}}.
 \label{eq:gap}
\end{equation}
This bounds $|(\gamma_i-\eta_j)-(c_A-c_B)|$ for individual ordinates
in the two prescribed windows.

\emph{Proof.} Write $P=p_{ij}\in[P_0,16]$. For $v\in\TT$, telescoping gives
$|v^j-v^l|\le |j-l|\,|1-v|$. Therefore
\[
 1-K_N(v)=\frac1{2N^2}\sum_{j,l=0}^{N-1}|v^j-v^l|^2
       \le\frac{N^2-1}{12}|1-v|^2 .
\]
Direct subtraction of the two Cayley fractions gives the exact cancellation
\[
 |1-z_i\bar v_j|^2=|z_i-v_j|^2=\frac{P}{d^2}|x_i-y_j|^2 .
\]
By positivity, $Q\ge P K_N(z_i\bar v_j)$. Combining the inequalities yields
$|x_i-y_j|^2\ge 12d^2(P-Q)/[(N^2-1)P^2]$. The derivative of
$(P-Q)/P^2$ has sign $\op{sgn}(2Q-P)$, so its minimum on the given
interval occurs at an endpoint. This proves \eqref{eq:gap}. \hfill$\square$

\subt{Verified applications}
Each row uses $c_A=c_B=0$ and $k=2$. The column headed $F/W^{2r}$
gives a rounded-up rational bound. The last column rounds \eqref{eq:gap} down.
\begin{center}\small
\begin{tabular}{@{}lrrrrrrr@{}}\toprule
Functions & $T$ & $a$ & $N$ & $r$ & degree & $F/W^{2r}<$ & gap $>$\\\midrule
$\zeta,\beta$ &20&14.5&50&5&250&0.682658&0.477\\
$\zeta,\beta$ &22&24&90&5&450&0.737047&0.361\\
$\zeta$ (distinct labels)&30&24&30&3&90&1.151405&1.29\\\bottomrule
\end{tabular}
\end{center}
A multiple zero would give a zero gap between distinct labels in the last row.
Hence all nontrivial zeta zeros with $|\gamma|\le30$ are simple. The method
also certifies their spacing without a list of ordinates.

\subt{A prescribed gap tested with three coefficients per center}
Take two copies of $\zeta$, $c_A=21$, $c_B=14$, $a=T=0.15$, $N=3$ and $r=k=1$.
Complex Taylor expansions through degree three give $F/W^2<0.761429$.
By \eqref{eq:gap}, every pair with $|\gamma-21|,|\eta-14|\le0.15$ satisfies
$|(\gamma-\eta)-7|>0.082$. The counts and sign checks on pages 6--7
place the first two positive zeta zeros in these windows.

\newpage
\HH{5. Arithmetic evaluation and an upper bound}
The verifiers use arbitrary-precision real and complex balls, with integer
arithmetic for the kernel coefficients \cite{arb}. For $\beta$, let
$\zeta(s,v)$ be the Hurwitz zeta function. We evaluate
\[
 H_\beta(s)=(4/\pi)^{(s+1)/2}\Gamma((s+1)/2)\,
             4^{-s}\{\zeta(s,1/4)-\zeta(s,3/4)\}.
\]
This is the completed $L$-function of the odd primitive character modulo
four. For $\zeta$, use $H_\zeta=\xi$. Both have real even restrictions to
$s=1/2+it$ as functions of real $t$.

For the certificates on page 5 the program computes $n[u^n]\log H(S(u))$ by
power-series arithmetic, then forms
\[
 b_{r,k}=-\sum_{j=-r}^r (-1)^j \binom{2r}{r+j}\ell_{k+j}.
\]
It forms the integer polynomial
$\bigl(\sum_{j=1-N}^{N-1}(N-|j|)X^{j+N-1}\bigr)^r$,
shifts its indices back, and divides by $N^{2r}$. Formula \eqref{eq:F}
and, where needed, the subtraction $b_{2r,0}$ now use only ball
operations. The code certifies both endpoint lower bounds in \eqref{eq:gap}
independently. The three rows use 2048, 3072 and 768 bits respectively. The shifted three-coefficient test uses 256 bits.
For the middle row, the two endpoint estimates begin with $0.3616249552$
and $0.4021952302$. The smaller one controls a guarantee uniform over
all the pairs in the window.

For an upper bound on the true minimum, the same program certifies the
following signs. The intervals enclose a sign change of each completion on
the critical line, so continuity gives one zero inside each interval.
\begin{center}
\begin{tabular}{@{}lccr@{}}\toprule
Completion & left endpoint & right endpoint & signs (left, right)\\\midrule
$\xi$ & 21.022 & 21.023 & $(-,+)$\\
$H_\beta$ & 21.450 & 21.451 & $(-,+)$\\\bottomrule
\end{tabular}
\end{center}
For these two zeros, $0<\eta-\gamma<21.451-21.022=0.429$.
This calculation also establishes that the minimum on page 1 ranges over a
nonempty finite set. The endpoints are rational sample heights. The proof
uses the function values at those heights. Exploratory approximate ordinates
helped select the filter parameters; the lower-bound verifier takes no
ordinates as inputs.

\subt{Reproducing the results}
With Python and \code{python-flint==0.9.0} installed, run
\code{python verify\_collisions.py} and \code{python verify\_zeta\_moments.py}.
Run Python with its ordinary assertion checks enabled. Both programs exit
with an error if a required interval inequality fails. Their complete output
is included. The general results on pages 2--5 are proved on paper; the
calculations certify the displayed specializations. The upstream Lean
project was inspected at source level. The reported upstream kernel checks are those
of the cited project and verification sources.

\newpage
\HH{6. How many zeta zeros are in the example?}
A small independent certificate counts the low zeta zeros using ten Taylor
coefficients. Write $\gamma_j$ for the positive ordinates in increasing
order with repetitions, and set $y_j=200/(1/4+\gamma_j^2)$ and
$M_k=\sum_j y_j^k$. Put $s(u)=(1-\sqrt{1-4u})/2$. The even Hadamard
product about $1/2$, divided by its value at $s=0$, gives
\[
 2\xi(s)=\prod_j\Big(1-\frac{s(1-s)}{1/4+\gamma_j^2}\Big),
 \qquad M_k=-k\,200^k[u^k]\log(2\xi(s(u))).
\]
Indeed the two factors for $\pm\gamma_j$ pair, and evenness removes the
linear exponential in the centered product. Absolute convergence follows
from $\sum_j\gamma_j^{-2}<\infty$. These inverse-power constructions have
a classical superzeta background \cite{voros}. At $s=0$ the code evaluates
$\xi(s)$ as $(s-1)\Gamma(1+s/2)\pi^{-s/2}\zeta(s)$.

\subt{Six coefficients give exactly one positive ordinate through 20}
Let $h(t)=200/(1/4+t^2)$, $L=0.9973$,
$U(y)=y^2-3y/8+1/100$, $g(y)=y^2(y-1/7)^2$.
Expanding the polynomials and substituting $M_1,\ldots,M_6$ gives
\[
 \sum_j y_j(y_j-L)U(y_j)^2>0.0000208,\qquad
 \sum_jg(y_j)-g(L)-g(h(20))<-0.0003234 .
\]
The first inequality forces $y_1>L>h(20)$. If also $y_2\ge h(20)$, the
first two terms of the second sum would already be at least
$g(L)+g(h(20))$, since $g\ge0$ and $g$ increases on $[1/7,\infty)$.
Thus exactly one label has $\gamma_j\le20$.

\subt{Ten coefficients give exactly two through 24}
For $0\le i,j\le4$ form $A_{ij}(t)=M_{i+j+2}-h(t)M_{i+j+1}$.
Its quadratic form at a polynomial $p$ of degree at most four is
$\sum_j y_j(y_j-h(t))p(y_j)^2$. It is nonpositive on the subspace
vanishing at the distinct $y_j>h(t)$, so its positive inertia gives a lower
bound on their number. All leading minors are nonzero. Including the empty
minor, their signs at $t=14.13474$ and $t=21.11$ are, respectively,
$(+,-,+,-,+,+)$ and $(+,-,-,+,-,-)$. The successive minor ratios are the
pivots of an $LDL^{\mathsf T}$ factorization. The counts of positive pivots
are one and two. Hence $y_1>h(14.13474)$ and $y_2>h(21.11)$. A third
label with $\gamma_j\le24$ would require at least $h(24)^6>0.0017478$,
whereas $M_6-h(14.13474)^6-h(21.11)^6<0.0016507$. Finally
$y_1\le M_{10}^{1/10}$ gives $\gamma_1>14.1344$. Thus the two zeros
are simple, one in $(14.1344,14.13474)$ and one in $(20,21.11)$.
Their known locations serve here as a check on finite-coefficient reasoning.

\newpage
\HH{7. Relation to prior work}
Li's zero sums, their extension to multisets by Bombieri and Lagarias, and
the Fourier description for automorphic $L$-functions by Lagarias already
connect coefficients to zero positions \cite{bl,lagarias}. The ordinary
Wiener lemma recovers the squared atomic masses of a measure from means
of its Fourier coefficients \cite{wiener}. Brown used finite Li data for
zero-free regions, and Ernvall-Hyt\"onen et al. studied generalized Li coefficients
for products of shifted functions \cite{brown,modified}. Suzuki developed related
positivity and moment identities for the zeta screw function \cite{suzuki}.
A particularly close preprint dated 24 November 2026 already gives a
finite-height simplicity criterion for $\zeta$, recovers the first multiple
ordinate from a trace defect, and reports a computation through height
$97.632$ \cite{trace}.

The specific contribution proposed here is formula \eqref{eq:F} for two
independently centered completions, with its exact finite-threshold
stabilization, convergence along all jointly diverging schedules, and
quantitative bound \eqref{eq:gap} on the distance from a prescribed gap.
The literature search found no earlier statement of this combination.
Priority remains open to further bibliographic checking. The general
simplicity and global disjointness questions remain open. A natural next
problem is to minimize the largest Taylor index needed for a specified
gap bound, by optimizing the nonnegative trigonometric polynomial.

\begingroup\small
\begin{thebibliography}{10}\setlength{\itemsep}{2pt}\setlength{\parskip}{0pt}
\bibitem{lean} \code{lsernau/grh-lean}, \emph{GRH for Dirichlet $L$-functions, in Lean 4},
readable edition, commit \code{72fead102f66} (25 Nov 2026).
\href{https://github.com/lsernau/grh-lean/blob/72fead102f663d5c7a407964a13dc0cf9d977207/grh/src/Statements.lean}{Statements},
\href{https://github.com/lsernau/grh-lean/blob/72fead102f663d5c7a407964a13dc0cf9d977207/grh/src/Main.lean}{proof entry points},
\href{https://github.com/lsernau/grh-lean/blob/72fead102f663d5c7a407964a13dc0cf9d977207/grh/GUIDE.md}{reading guide}.
See A. Sutherland's \href{https://math.mit.edu/~drew/shortgrh/}{independent checks of both editions} and the \href{https://leanprover-community.github.io/blog/posts/riemann-hypothesis/}{Lean FRO and Mathlib statement} on Alp\"oge's original proof.
\bibitem{wiki} Wikipedia, \href{https://en.wikipedia.org/wiki/Riemann_hypothesis}{\emph{Riemann hypothesis}}, sections on the formal proof and equivalent analytic criteria, consulted 26 Nov 2026.
\bibitem{bl} E. Bombieri and J. C. Lagarias, \href{https://doi.org/10.1006/jnth.1999.2392}{Complements to Li's criterion for the Riemann hypothesis}, J. Number Theory 77 (1999), 274--287.
\bibitem{lagarias} J. C. Lagarias, \href{https://doi.org/10.5802/aif.2311}{Li coefficients for automorphic $L$-functions}, Ann. Inst. Fourier 57 (2007), 1689--1740, especially Theorem 2.1.
\bibitem{sek} S. K. Sekatskii, \href{https://umj.imath.kiev.ua/index.php/umj/article/view/2139}{Generalized Bombieri--Lagarias' theorem and generalized Li's criterion with its arithmetic interpretation}, Ukrainian Math. J. 66 (2014).
\bibitem{wiener} C. Cuny, T. Eisner and B. Farkas, \href{https://imacm.uni-wuppertal.de/fileadmin/imacm/preprints/2018/imacm_18_22.pdf}{Wiener's lemma along primes and other subsequences}, Theorem 1.1.
\bibitem{suzuki} M. Suzuki, \href{https://doi.org/10.1112/jlms.12785}{Aspects of the screw function corresponding to the Riemann zeta-function}, J. London Math. Soc. (2023).
\bibitem{trace} \href{https://hexagonmath.org/pdf/2611.00996v1}{A finite arithmetic test for simplicity of the Riemann zeros, after Suzuki}, preprint (24 Nov 2026), Theorem 1.
\bibitem{arb} F. Johansson et al., \href{https://python-flint.readthedocs.io/en/latest/arb_series.html}{python-flint, real ball power series}, version 0.9.0.
\bibitem{brown} F. C. S. Brown, \href{https://doi.org/10.1016/j.jnt.2004.07.016}{Li's criterion and zero-free regions of $L$-functions}, J. Number Theory 111 (2005), 1--32.
\bibitem{modified} A.-M. Ernvall-Hyt\"onen, A. Od\v{z}ak, L. Smajlovi\'c and M. Su\v{s}i\'c, \href{https://arxiv.org/abs/1410.4384}{On the modified Li criterion for a certain class of $L$-functions}, J. Number Theory 156 (2015).
\bibitem{voros} A. Voros, \href{https://doi.org/10.5802/aif.1955}{Zeta functions for the Riemann zeros}, Ann. Inst. Fourier 53 (2003), 665--699.
\end{thebibliography}
\endgroup
\end{document}
